\documentclass[../../main.tex]{subfiles}
\graphicspath{{../../image/}{../../figure/}}
\begin{document}

\section{Hugoniot 轨迹}\label{section:ch08-hugoniot轨迹}

状态空间中可由激波连接到固定点 $\bar U$ 的点 $U$ 所成的集合，称为 $\bar U$ 的 \emph{Hugoniot 轨迹}。在系统任一严格双曲性点 $U$ 的邻域内，它具有简单的几何结构。

\begin{theorem}[激波曲线]\label{theorem:ch08-激波曲线}
给定状态 $\bar U\in\mathcal O$，设特征速度 $\lambda_i(\bar U)$ 是 $\mathrm{D}F(\bar U)$ 的单特征值。则在状态空间中存在一条 $C^3$ 曲线 $U=W_i(\tau)$（称为过 $\bar U$ 的 $i$-激波曲线），以及一个 $C^2$ 函数 $s=s_i(\tau)$，二者均定义在 $0$ 的某个邻域上，并具有下述性质：状态 $U$ 与 $\bar U$ 可由一个中等强度、速度为 $s$ 的 $i$-激波相连，当且仅当存在某个 $\tau$ 使 $U=W_i(\tau)$、$s=s_i(\tau)$。进而 $W_i(0)=\bar U$，且
\begin{align}
s_i(0)&=\lambda_i(\bar U),
\label{align:ch08-si0}\\
\dot s_i(0)&=\tfrac12\mathrm{D}\lambda_i(\bar U)R_i(\bar U),
\label{align:ch08-sdot}\\
\dot W_i(0)&=R_i(\bar U),
\label{align:ch08-wdot}\\
\ddot W_i(0)&=\mathrm{D}R_i(\bar U)R_i(\bar U).
\label{align:ch08-wddot}
\end{align}
当需要标明该激波曲线的原点时，将使用更明确的记号 $W_i(\tau;\bar U)$、$s_i(\tau;\bar U)$。
\end{theorem}

\begin{proof}
回顾\hyperref[section:ch08-强激波弱激波中等强度激波]{8.1 节}中发展的记号，特别是\eqref{align:ch08-s中等}与\eqref{align:ch08-跳跃方向}。状态 $U$ 与 $\bar U$ 可由速度为 $s$ 的 $i$-激波相连，当且仅当
\begin{align}
U&=\bar U+\tau S_i(\bar U,U),
\label{align:ch08-隐式1}\\
s&=\mu_i(\bar U,U).
\label{align:ch08-隐式2}
\end{align}
据此考虑定义在 $\mathcal O\times\mathbb R$ 上的函数
\begin{align}
H(U,\tau)=U-\bar U-\tau S_i(\bar U,U),
\label{align:ch08-H}
\end{align}
注意 $H(\bar U,0)=0$，$\mathrm{D}H(\bar U,0)=I$。于是由隐函数定理，存在状态空间中的曲线 $U=W_i(\tau)$，$W_i(0)=\bar U$，使得对 $\tau$ 充分接近 $0$，$H(U,\tau)=0$ 当且仅当 $U=W_i(\tau)$。我们再定义
\begin{align}
s_i(\tau)=\mu_i(\bar U,W_i(\tau)).
\label{align:ch08-si定义}
\end{align}
特别地，$s_i(0)=\mu_i(\bar U,\bar U)=\lambda_i(\bar U)$。进而，将\eqref{align:ch08-隐式1} 关于 $\tau$ 求导并置 $\tau=0$，得 $\dot W_i(0)=S_i(\bar U,\bar U)=R_i(\bar U)$。为建立剩余的\eqref{align:ch08-sdot} 与\eqref{align:ch08-wddot}，借助\eqref{align:ch08-mu展开} 与\eqref{align:ch08-S展开}，得
\begin{align}
s_i(\tau)&=\lambda_i\!\left(\tfrac12(\bar U+W_i(\tau))\right)+O(\tau^2)
=\lambda_i(\bar U)+\tfrac12\tau\mathrm{D}\lambda_i(\bar U)R_i(\bar U)+O(\tau^2),
\label{align:ch08-si展开}\\
W_i(\tau)&=\bar U+\tau R_i\!\left(\tfrac12(\bar U+W_i(\tau))\right)+O(\tau^3)
=\bar U+\tau R_i(\bar U)+\tfrac12\tau^2\mathrm{D}R_i(\bar U)R_i(\bar U)+O(\tau^3).
\label{align:ch08-wi展开}
\end{align}
证明完毕。
\end{proof}

特别地，若 $\bar U$ 是系统\eqref{align:ch08-守恒律一维} 的一个严格双曲性点，则\cref{theorem:ch08-激波曲线} 蕴涵 $\bar U$ 的 Hugoniot 轨迹是 $n$ 条激波曲线之并，每个特征族对应一条。

上述构造的激波曲线一般局限于中等强度激波的范围，这是由于使用了隐函数定理——它仅在以 $|\tau|$ 度量的激波强度充分小时适用：$|\tau|<\delta$，其中 $\delta$ 依赖于 $S_i$ 的 $C^1$ 范数，而后者又可由 $\mathrm{D}F$ 的 $C^1$ 范数以及 $\lambda_i$ 与其他特征速度间距的倒数来估计。尽管如此，在特殊系统中，人们往往可借助精细的解析或拓扑论证，或显式计算，将激波曲线延拓到强激波范围。例如，对系统\eqref{align:ch07-p系统}，将\eqref{align:ch08-弹性rh} 与\eqref{align:ch08-弹性s} 联立，可推知状态空间中任一点 $(\bar u,\bar v)$ 的 Hugoniot 轨迹由两条曲线
\begin{align}
v=\bar v\pm\sqrt{\big[\sigma(u)-\sigma(\bar u)\big](u-\bar u)},
\label{align:ch08-弹性轨迹}
\end{align}
组成，定义在 $u$ 的整个取值范围上。

另一个值得注意的、激波曲线可延拓到强激波领域的系统是\eqref{align:ch07-热弹性}。回顾\hyperref[section:ch08-强激波弱激波中等强度激波]{8.1 节}末的讨论，我们推知过状态 $(\bar u,\bar v,\bar s)$ 的 2-激波曲线由 Rankine--Hugoniot 跳跃条件确定：
\begin{align}
v=\bar v,\qquad p(u,s)=p(\bar u,\bar s).
\label{align:ch08-2激波曲线}
\end{align}
由于 $p_u<0$，\eqref{align:ch08-2激波曲线} 描述了一条以 $s$ 参数化的简单曲线。至于另外两条激波曲线，Hugoniot 方程\eqref{align:ch08-hugoniot方程}（取 $(u_-,s_-)=(\bar u,\bar s)$、$(u_+,s_+)=(u,s)$）确定了 Hugoniot 轨迹在 $u$-$s$ 平面上的投影。将此投影实现为曲线 $s=s(u)$，关于 $u$ 求导并利用\eqref{align:ch07-热力学}，得
\begin{align}
\big[2\theta-\theta_u(u-\bar u)\big]\frac{\mathrm{d}s}{\mathrm{d}u}=p-\bar p-p_u(u-\bar u).
\label{align:ch08-sde}
\end{align}
由于 $\theta>0$，\eqref{align:ch08-sde} 在 $\bar u$ 的某个邻域上诱导出一条简单曲线 $s=s(u)$。当状态方程在其定义域上满足 $2\theta-u\theta_u>0$ 时（理想气体\eqref{align:ch02-2-5-20} 即是如此），曲线 $s(u)$ 延拓到强激波范围。然而，若允许 $2\theta-u\theta_u$ 变号，则在遇到 $2\theta-\theta_u(u-\bar u)$ 与 $p-\bar p-p_u(u-\bar u)$ 同时为零的奇点时，$s(u)$ 可能分裂成无穷多支。1-激波与 3-激波曲线的速度分量由\eqref{align:ch08-气体rh} 给出：
\begin{align}
v(u)=\bar v\pm\sqrt{-\big[p(u,s(u))-p(\bar u,\bar s)\big](u-\bar u)}.
\label{align:ch08-vu}
\end{align}

正如我们将在\hyperref[section:ch08-熵激波容许性准则]{8.5 节}看到的，以 $u$ 参数化激波曲线将阐明强激波的容许性。为给该研究铺垫，我们对\eqref{align:ch08-sde} 求导，得
\begin{align}
\big[2\theta-\theta_u(u-\bar u)\big]\frac{\mathrm{d}^2s}{\mathrm{d}u^2}
+\big[2\theta_s-\theta_{us}(u-\bar u)\big]\left(\frac{\mathrm{d}s}{\mathrm{d}u}\right)^2
+2\big[\theta_u-\theta_{uu}(u-\bar u)\big]\frac{\mathrm{d}s}{\mathrm{d}u}
=-p_{uu}(u-\bar u).
\label{align:ch08-sde2}
\end{align}
以 $u$ 参数化的激波速度 $r=r(u)$ 由\eqref{align:ch08-气体r} 确定，其中取 $(u_-,s_-)=(\bar u,\bar s)$、$(u_+,s_+)=(u,s(u))$。利用\eqref{align:ch08-sde}，推得
\begin{align}
r\frac{\mathrm{d}r}{\mathrm{d}u}=\frac{\theta}{(u-\bar u)^2}\frac{\mathrm{d}s}{\mathrm{d}u}.
\label{align:ch08-rdr}
\end{align}

回到一般系统\eqref{align:ch08-守恒律一维}，注意到上面引入的 $i$-激波曲线与\cref{theorem:ch07-稀疏波曲线存在唯一}中定义的 $i$-稀疏波曲线具有共同特征。确实，回顾\cref{theorem:ch07-稀疏波曲线存在唯一} 与\cref{theorem:ch08-激波曲线}，特别是将\cref{theorem:ch07-稀疏波曲线存在唯一} 与\eqref{align:ch08-wdot}、\eqref{align:ch08-wddot} 作比较，我们推得

\begin{theorem}[激波曲线与稀疏波曲线的二阶切触]\label{theorem:ch08-二阶切触}
设 $\bar U\in\mathcal O$ 是守恒律双曲系统\eqref{align:ch08-守恒律一维} 的 $i$-特征族真非线性点，且 $\lambda_i(\bar U)$ 是 $\mathrm{D}F(\bar U)$ 的单特征值。规范化 $R_i$，使\cref{theorem:ch07-稀疏波曲线存在唯一} 在 $\bar U$ 的某个邻域上成立。则由\cref{theorem:ch07-稀疏波曲线存在唯一} 定义的 $i$-稀疏波曲线 $V_i$ 与由\cref{theorem:ch08-激波曲线} 定义的 $i$-激波曲线 $W_i$，在 $\bar U$ 处具有二阶切触。
\end{theorem}

回顾由\cref{theorem:ch07-稀疏波曲线Riemann刻画}，$i$-Riemann 不变量沿 $i$-稀疏波曲线为常数。同时，如上所示，$i$-激波曲线与 $i$-稀疏波曲线非常接近。于是自然预期 $i$-Riemann 不变量沿 $i$-激波曲线变化非常缓慢。确实，

\begin{theorem}[弱激波跨跃中 Riemann 不变量的三阶小量]\label{theorem:ch08-riemann三阶}
任一 $i$-Riemann 不变量跨过弱 $i$-激波的跳跃，是激波强度的三阶小量。
\end{theorem}

\begin{proof}
设 $\lambda_i(\bar U)$ 是 $\mathrm{D}F(\bar U)$ 的单特征值，考虑过 $\bar U$ 的 $i$-激波曲线 $W_i$。对任一 $i$-Riemann 不变量 $w$，沿曲线 $W_i(\cdot)$ 求导，
\begin{align}
\dot w&=\mathrm{D}w\,\dot W_i,
\label{align:ch08-wdotw}\\
\ddot w&=\dot W_i^{\top}\mathrm{D}^2 w\,\dot W_i+\mathrm{D}w\,\ddot W_i.
\label{align:ch08-wddotw}
\end{align}
由\eqref{align:ch08-wdot} 与\eqref{align:ch07-Riemann不变量定义}，在 $\tau=0$ 处 $\dot w=0$。

现对\eqref{align:ch07-Riemann不变量定义} 作用 $\mathrm{D}$，再从右侧乘以 $R_i$，推得恒等式
\begin{align}
R_i^{\top}\mathrm{D}^2 w\,R_i+\mathrm{D}w\,\mathrm{D}R_i\,R_i=0.
\label{align:ch08-riemann恒等式}
\end{align}
联立\eqref{align:ch08-wddotw}、\eqref{align:ch08-wdot}、\eqref{align:ch08-wddot} 与\eqref{align:ch08-riemann恒等式}，得在 $\tau=0$ 处 $\ddot w=0$。证明完毕。
\end{proof}

在系统配有 Riemann 不变量坐标系 $(w_1,\cdots,w_n)$ 的特殊情形，我们可如下计算弱 $i$-激波（$i\neq j$）下 $w_j$ 跳跃的首项。Rankine--Hugoniot 条件读作
\begin{align}
F(W_i(\tau))-F(\bar U)=s_i(\tau)\big[W_i(\tau)-\bar U\big].
\label{align:ch08-rh曲线}
\end{align}
关于 $\tau$ 求导，得
\begin{align}
\big[\mathrm{D}F(W_i(\tau))-s_i(\tau)I\big]\dot W_i(\tau)=\dot s_i(\tau)\big[W_i(\tau)-\bar U\big].
\label{align:ch08-drh曲线}
\end{align}
从左侧乘以 $\mathrm{D}w_j(W_i)$，得
\begin{align}
(\lambda_j-s_i)\dot w_j=\dot s_i\,\mathrm{D}w_j\big[W_i-\bar U\big].
\label{align:ch08-drh2}
\end{align}
接着关于 $\tau$ 求导，得
\begin{align}
(\lambda_j-s_i)\ddot w_j+(\dot\lambda_j-2\dot s_i)\dot w_j
=\ddot s_i\,\mathrm{D}w_j\big[W_i-\bar U\big]
+\dot s_i\,\dot W_i^{\top}\mathrm{D}^2 w_j\big[W_i-\bar U\big].
\label{align:ch08-drh3}
\end{align}
关于 $\tau$ 再求导并置 $\tau=0$。利用\eqref{align:ch08-si0}、\eqref{align:ch08-sdot}、\eqref{align:ch08-wdot}、\eqref{align:ch07-链式法则}，以及由\cref{theorem:ch08-riemann三阶} 知 $\dot w_j$、$\ddot w_j$ 在 $0$ 处为零，推得
\begin{align}
\dddot w_j=\frac12\,\frac{1}{\lambda_j-\lambda_i}\frac{\partial\lambda_i}{\partial w_i}\,R_i^{\top}\mathrm{D}^2 w_j\,R_i,
\label{align:ch08-riemann三阶展开}
\end{align}
其中 $\dddot w_j$ 在 $0$ 处取值，右端在 $\bar U$ 处取值。

回到一般情形，我们接下来研究激波速度函数 $s_i(\tau)$ 沿 $i$-激波曲线如何演化。从左侧以 $L_i(W_i(\tau))$ 乘以\eqref{align:ch08-drh曲线}，得
\begin{align}
\big[\lambda_i(W_i(\tau))-s_i(\tau)\big]L_i(W_i(\tau))\dot W_i(\tau)
=\dot s_i(\tau)L_i(W_i(\tau))\big[W_i(\tau)-\bar U\big].
\label{align:ch08-drhli}
\end{align}
对充分接近 $0$ 但 $\tau\neq0$ 的 $\tau$，由\eqref{align:ch08-wdot} 知
\begin{align}
L_i(W_i(\tau))\dot W_i(\tau)>0,\qquad \tau L_i(W_i(\tau))\big[W_i(\tau)-\bar U\big]>0.
\label{align:ch08-符号}
\end{align}
在应用中，\eqref{align:ch08-符号} 往往对宽广范围的 $\tau$ 继续成立，常延拓到强激波范围。此时，\eqref{align:ch08-drhli} 与\eqref{align:ch08-drh曲线} 立即给出下述

\begin{lemma}[激波速度沿激波曲线的符号]\label{lemma:ch08-激波速度符号}
设\eqref{align:ch08-符号} 成立。则
\begin{align}
\dot s_i(\tau)>0\quad&\text{当且仅当}\quad \tau\big[\lambda_i(W_i(\tau))-s_i(\tau)\big]>0,
\label{align:ch08-sdotpos}\\
\dot s_i(\tau)=0\quad&\text{当且仅当}\quad \lambda_i(W_i(\tau))=s_i(\tau).
\label{align:ch08-sdotzero}
\end{align}
进而，$\dot s_i(\tau)=0$ 蕴涵 $\dot W_i(\tau)$ 与 $R_i(W_i(\tau))$ 共线。
\end{lemma}

为看清 $s_i$ 在 $\dot s_i$ 为零的点附近如何变化，将\eqref{align:ch08-drhli} 关于 $\tau$ 求导，并在任一使 $\dot s_i(\tau)=0$ 的 $\tau$ 处求值。由于此时 $s_i(\tau)=\lambda_i(W_i(\tau))$ 且 $\dot W_i(\tau)=aR_i(W_i(\tau))$，回顾\eqref{align:ch07-左右向量规范化}，推得
\begin{align}
\ddot s_i(\tau)L_i(W_i(\tau))\big[W_i(\tau)-\bar U\big]
=a^2\mathrm{D}\lambda_i(W_i(\tau))R_i(W_i(\tau)),
\label{align:ch08-sddot}
\end{align}
由此可知，在 $\dot s_i=0$ 的点处，$\ddot s_i$ 与 $\tau\,\mathrm{D}\lambda_i R_i$ 同号。

由\cref{lemma:ch08-激波速度符号}，$s_i$ 为常数蕴涵 $i$-激波曲线是向量场 $R_i$ 的一条积分曲线，沿它 $\lambda_i$ 为常数。因此，沿这样一条激波曲线的所有点都是 $i$-特征族的线性退化状态。该命题的逆命题也成立：

\begin{theorem}[线性退化族的激波曲线]\label{theorem:ch08-线性退化激波曲线}
设守恒律双曲系统\eqref{align:ch08-守恒律一维} 的 $i$-特征族线性退化，且 $\lambda_i(\bar U)$ 是 $\mathrm{D}F(\bar U)$ 的单特征值。则过 $\bar U$ 的 $i$-激波曲线 $W_i$ 就是过 $\bar U$ 的 $R_i$ 的积分曲线。事实上，在适当参数化下，$W_i$ 是微分方程
\begin{align}
\dot W_i=R_i(W_i)
\label{align:ch08-riemann不变常数}
\end{align}
满足初始条件 $W_i(0)=\bar U$ 的解。沿 $W_i$，特征速度 $\lambda_i$ 及所有 $i$-Riemann 不变量均为常数。激波速度函数 $s_i$ 也为常数：
\begin{align}
s_i(\tau)=\lambda_i(W_i(\tau))=\lambda_i(\bar U).
\label{align:ch08-si常数}
\end{align}
\end{theorem}

\begin{proof}
设 $W_i$ 是\eqref{align:ch08-riemann不变常数} 满足 $W_i(0)=\bar U$ 的解。则
\begin{align}
\big[\mathrm{D}F(W_i(\tau))-\lambda_i(W_i(\tau))I\big]\dot W_i(\tau)=0.
\label{align:ch08-wdot方程}
\end{align}
由于 $\mathrm{D}\lambda_i(U)R_i(U)=0$，$\dot\lambda_i(W_i(\tau))=0$，从而 $\lambda_i(W_i(\tau))=\lambda_i(\bar U)$。将\eqref{align:ch08-wdot方程} 从 $0$ 到 $\tau$ 积分，得
\begin{align}
F(W_i(\tau))-F(\bar U)=\lambda_i(\bar U)\big[W_i(\tau)-\bar U\big],
\label{align:ch08-积分rh}
\end{align}
这表明 $W_i$ 是过 $\bar U$ 的 $i$-激波曲线，相应激波速度函数 $s_i$ 由\eqref{align:ch08-si常数} 给出。证明完毕。
\end{proof}

\cref{theorem:ch08-线性退化激波曲线} 的下述重要推论给出了线性退化的另一种刻画：

\begin{corollary}[线性退化族的行波解]\label{corollary:ch08-行波解}
当双曲系统\eqref{align:ch08-守恒律一维} 的 $i$-特征族线性退化时，对 $i$-特征速度 $\lambda_i$ 取值范围内的任一 $\sigma$，存在行波解
\begin{align}
U(x,t)=V(x-\sigma t).
\label{align:ch08-行波}
\end{align}
\end{corollary}

\begin{proof}
设对某状态 $\bar U$ 有 $\sigma=\lambda_i(\bar U)$。考虑过 $\bar U$、满足\eqref{align:ch08-riemann不变常数} 的 $i$-激波曲线 $W_i$。任取 $C^1$ 函数 $\tau=\tau(\xi)$，并以 $V(\xi)=W_i(\tau(\xi))$ 通过\eqref{align:ch08-行波} 定义 $U$。由\eqref{align:ch08-riemann不变常数} 与\eqref{align:ch08-si常数}，
\begin{align}
\partial_t U+\partial_x F(U)=\frac{\mathrm{d}\tau}{\mathrm{d}\xi}
\big[\mathrm{D}F(W_i(\tau))-\lambda_i(W_i(\tau))I\big]R_i(W_i(\tau))=0.
\label{align:ch08-行波代入}
\end{align}
证明完毕。
\end{proof}

自然要问：在缺乏线性退化时，$i$-激波曲线是否可能是向量场 $R_i$ 的积分曲线？结果表明，这只有在非常特殊的情形下才可能发生：

\begin{theorem}[激波曲线与积分曲线重合的条件]\label{theorem:ch08-重合条件}
对双曲系统\eqref{align:ch08-守恒律一维}，设 $\bar U$ 是 $i$-特征族的真非线性状态，且 $\lambda_i(\bar U)$ 是 $\mathrm{D}F(\bar U)$ 的单特征值。过 $\bar U$ 的 $i$-激波曲线与向量场 $R_i$ 过 $\bar U$ 的积分曲线（即 $i$-稀疏波曲线）重合，当且仅当后者在状态空间中是一条直线。
\end{theorem}

\begin{proof}
若过 $\bar U$ 的 $i$-激波曲线 $W_i$ 与 $R_i$ 过 $\bar U$ 的积分曲线重合，则 $\dot W_i(\tau)$ 必须与 $R_i(W_i(\tau))$ 共线。此时，\eqref{align:ch08-drh曲线} 蕴涵
\begin{align}
\big[\lambda_i(W_i(\tau))-s_i(\tau)\big]\dot W_i(\tau)
=\dot s_i(\tau)\big[W_i(\tau)-\bar U\big].
\label{align:ch08-共线}
\end{align}
对 $\tau$ 接近 $0$ 但 $\tau\neq0$，由真非线性性有 $\lambda_i(W_i(\tau))\neq s_i(\tau)$。因此，\eqref{align:ch08-共线} 蕴涵 $W_i$ 的图像是一条过 $\bar U$ 的直线。

反之，设 $R_i$ 过 $\bar U$ 的积分曲线是一条直线，可参数化为 $U=W_i(\tau)$，其中 $W_i$ 是某个满足 $W_i(0)=\bar U$，以及\eqref{align:ch08-wdot} 与\eqref{align:ch08-wddot} 的光滑函数（注意 $\mathrm{D}R_i(\bar U)R_i(\bar U)$ 必然与 $R_i(\bar U)$ 共线）。于是我们可确定一个标量值函数 $s_i(\tau)$，使得
\begin{align}
F(W_i(\tau))-F(\bar U)
&=\int_0^\tau \mathrm{D}F(W_i(\zeta))\dot W_i(\zeta)\,\mathrm{d}\zeta
=\int_0^\tau \lambda_i(W_i(\zeta))\dot W_i(\zeta)\,\mathrm{d}\zeta
=s_i(\tau)\big[W_i(\tau)-\bar U\big].
\label{align:ch08-直线积分}
\end{align}
从而 $W_i$ 是过 $\bar U$ 的 $i$-激波曲线。证明完毕。
\end{proof}

尽管特殊，每个特征族的激波曲线与稀疏波曲线重合的双曲守恒律系统类，包含一些值得注意的例子。例如，考虑电渗系统\eqref{align:ch07-电泳}。注意，对 $i=1,\cdots,n$，由\eqref{align:ch07-电泳wi方程} 或\eqref{align:ch07-电泳wn} 确定的 $i$-Riemann 不变量 $W_i$ 的等值面是超平面。特别地，对 $i=1,\cdots,n$，向量场 $R_i$ 的积分曲线是由 $n-1$ 个 Riemann 不变量 $w_1,\cdots,w_{i-1},w_{i+1},\cdots,w_n$ 的等值超平面相交而成的直线。因此，\cref{theorem:ch08-重合条件} 的条件适用于系统\eqref{align:ch07-电泳}。

在存在多重特征速度时，Hugoniot 轨迹可能包含多维流形，以代替激波曲线。与此相关，考虑模型系统\eqref{align:ch07-脐点模型} 是有启发意义的，其原点是一个脐点。当状态 $(\bar u,\bar v)$ 与状态 $(u,v)$ 由速度为 $s$ 的激波相连时，Rankine--Hugoniot 条件读作
\begin{align}
\begin{cases}
(u^2+v^2)u-(\bar u^2+\bar v^2)\bar u=s(u-\bar u)\\
(u^2+v^2)v-(\bar u^2+\bar v^2)\bar v=s(v-\bar v).
\end{cases}
\label{align:ch08-umiblic}
\end{align}
注意，当 $(\bar u,\bar v)\neq(0,0)$ 时，$(\bar u,\bar v)$ 的 Hugoniot 轨迹由圆周 $u^2+v^2=\bar u^2+\bar v^2$（沿它激波速度为常数 $s=\bar u^2+\bar v^2$）以及直线 $\bar v u=\bar u v$（它将 $(\bar u,\bar v)$ 连到原点）组成。于是，1-特征族提供了一个适用\cref{theorem:ch08-线性退化激波曲线} 的情形，同时 2-特征族满足\cref{theorem:ch08-重合条件} 的假设。另一方面，脐点 $(0,0)$ 的 Hugoniot 轨迹是整个平面，因为任一点 $(u,v)$ 都可由速度 $s=u^2+v^2$ 的激波连到 $(0,0)$。

并非所有严格双曲性失效的系统都表现出同样行为。例如，原点处严格双曲性也失效的系统
\begin{align}
\begin{cases}
\partial_t u+\partial_x\big[2(u^2+v^2)u\big]=0\\
\partial_t v+\partial_x\big[(u^2+v^2)v\big]=0,
\end{cases}
\label{align:ch08-非脐点}
\end{align}
其 $(0,0)$ 的 Hugoniot 轨迹由两条直线组成，即 $u$ 轴与 $v$ 轴。

\end{document}
